\chapter{Simulation, tests and importance of effects}\label{ch:tools}

This chapter describes the commands that describe the data before a fit and work with a fitted model or simulate one: the description of the data, holding effects fixed, simulating a model at chosen parameters, the score-type test for effects that were left out of a model, the test for differences between periods, and the relative importance of the effects. All of them follow the corresponding procedures of RSiena; the comparison is in Chapter~\ref{ch:validation}.

\section{How the commands are organised}

A property of the data is detected from the data: a two-mode network or a non-directed network is recognised from the way the network was declared and needs no option. An option changes how a model is specified or fitted (\opt{behavior()}, \opt{present()}, \opt{fix()}). A different job is a subcommand (\texttt{nwsaom multinet}, \texttt{nwsaom simulate}) or a postestimation command (\texttt{estat}). \texttt{nwsaom multiplex} is deprecated and forwards to \texttt{nwsaom multinet}.

\section{Describing the data: \texttt{nwsaom describe}}\label{sec:describe}

Before a fit, \texttt{nwsaom describe} prints what RSiena prints before an estimation: for every wave the number of ties, the density and the average degree; for every period the numbers of ties that are kept, created and dissolved, the Jaccard index and the Hamming distance; and, with \opt{behavior()}, the mean and standard deviation of the behavior in every wave and the numbers of actors whose behavior goes up, goes down or stays the same.
\begin{lstlisting}
. nwwebuse glasgow, nwclear
. nwsaom describe, waves(glasgow1 glasgow2 glasgow3) behavior(alcohol1 alcohol2 alcohol3)
Network data: 3 waves, 50 actors, directed
-----------------------------------------------------
        wave |     ties     density   avg. degree
-------------+---------------------------------------
           1 |      113      0.0461         2.260
           2 |      116      0.0473         2.320
           3 |      122      0.0498         2.440
-----------------------------------------------------

----------------------------------------------------------------------
      period |    kept   created  dissolved   Jaccard   Hamming
-------------+--------------------------------------------------------
        1->2 |      57        59         56     0.331       115
        2->3 |      66        56         50     0.384       106
----------------------------------------------------------------------

--------------------------------------------------
    behavior |        mean          sd     missing
-------------+------------------------------------
      wave 1 |       2.880       1.118           0
      wave 2 |       3.100       1.199           0
      wave 3 |       3.360       0.985           0
--------------------------------------------------

--------------------------------------------------
      change |          up        down        same
-------------+------------------------------------
        1->2 |          16           6          28
        2->3 |          18           9          23
--------------------------------------------------
\end{lstlisting}
A period with a Jaccard index below about 0.3 has few stable ties and carries little information about the change process; here the two periods are at 0.33 and 0.38. The results are in \texttt{r(periods)}, \texttt{r(waves)} and \texttt{r(behchange)}. For a two-mode network the ties are the attendances (the cells of the actors-by-events matrix), the density is taken over those cells and the average degree is the mean number of events per actor; the option \opt{actormode(\#)} names the mode of the actors (default 2, as in \texttt{nwsaom}), and \texttt{r(n\_actors)} and \texttt{r(n\_events)} hold the numbers of actors and events. The description of a stacked data set with groups (Chapter~\ref{ch:groups}) is that of the combined network; describe the groups one by one by making a network of each.

\section{Holding effects fixed: \opt{fix()}}

\opt{fix(\emph{effects})} keeps the named effects at their \opt{theta0()} values. They are not estimated and have no standard error; this is RSiena's \texttt{fix = TRUE}. \opt{fix(\_all)} fixes every effect. The names are those of the effects in the call; the effects of a co-evolving behavior are named by their token (\opt{fix(avalt)}) or by the label in the coefficient table (\texttt{beh\_avalt}).
\begin{lstlisting}
. nwsaom, wave1(glasgow1) wave2(glasgow2) outdegree reciprocity transtrip theta0(-2 1 0) fix(reciprocity) k3(300) seed(1)
   outdegree |  -2.311036   .1114178   -20.74   0.000
 reciprocity |          1          .        .       .
   transtrip |   .7747333   .0996184     7.78   0.000
(output omitted)
Overall maximum convergence ratio: 0.188 (criterion: below 0.25; met)
Convergence t-ratios (criterion: all below 0.1 in absolute value): not met for outdegree (0.132) reciprocity (-1.656) transtrip (0.185)
Convergence is good only if both criteria are met. Refit starting from these estimates with theta0(); see nwsaom.
\end{lstlisting}
The fixed coefficient equals its start value and its variance is 0. The overall maximum convergence ratio leaves out the fixed parameter, but the line of the $t$-ratios checks the $t$-ratios of all parameters, the fixed one included, so it lists \texttt{reciprocity} ($-1.656$); a large $t$-ratio of a fixed parameter is expected and is what the score test below examines. The $t$-ratios of the two free parameters, \texttt{outdegree} (0.132) and \texttt{transtrip} (0.185), are above 0.1 as well, so this fit has not converged by the two criteria; before it is reported it would be restarted from its estimates with \opt{theta0()}, as Chapter~\ref{ch:estimation} explains. Fixing is used for the score test below, for fitting a model with a known part, and for the simulation-only mode.

\section{Simulation: \texttt{estat simulate} and \texttt{nwsaom simulate}}\label{sec:sim}

\texttt{estat simulate} simulates the fitted model forward. Every period starts at the observed start wave of the period and runs for the period's length at the fitted coefficients and rates, the same simulation that \texttt{estat gof} uses; the default is 100 replications.
\begin{lstlisting}
. nwsaom, wave1(glasgow1) wave2(glasgow2) outdegree reciprocity transtrip k3(200) seed(1)
. estat simulate, nsim(20) seed(5)
Simulation from the fitted model: 20 replication(s) of 1 period(s)

    statistic |       Sim       Obs
--------------+---------------------
Period 1      |
         Ties |     112.4       116
       Mutual |      34.0        35
   Transitive |      77.8        88
------------------------------------
Sim is the mean over the replications; Obs is the observed end wave.
\end{lstlisting}
Options: \opt{theta(\emph{matname})} replaces the coefficients (a row vector in the order of \texttt{e(b)}), \opt{rates()} and \opt{ratesbeh()} replace the rates, \opt{generate(\emph{stub})} saves the simulated networks as networks named \emph{stub}\texttt{\_}\emph{s}\texttt{\_p}\emph{period} (and the behaviors as variables of the same names). The result \texttt{r(sim)} has one row per replication and period (number of ties, mutual dyads, transitive triples, mean behavior), \texttt{r(obs)} the same for the observed end waves.

\paragraph{Simulation-only mode.} \texttt{nwsaom simulate} simulates at coefficients you choose, without any estimation (RSiena's \texttt{siena07} with \texttt{simOnly = TRUE}). The model is specified exactly as in \texttt{nwsaom}; \opt{theta()} gives one coefficient per network effect in the order of the coefficient table of a fit of the same model (the order in which \nwsaom{} lists the effects, which is not always the order in which you typed them), \opt{behtheta()} one per behavior effect, \opt{rates()} the rate of each period.
\begin{lstlisting}
. nwsaom simulate, wave1(glasgow1) wave2(glasgow2) outdegree reciprocity transtrip theta(-2.3 2 0.5) rates(5) nsim(30) seed(3)
    statistic |       Sim       Obs
--------------+---------------------
Period 1      |
         Ties |     127.1       116
       Mutual |      36.3        35
   Transitive |      72.0        88
------------------------------------
Sim is the mean over the replications; Obs is the observed end wave.
\end{lstlisting}
Internally the model is set up by a short call of \texttt{nwsaom} with all coefficients fixed, so every data option of \texttt{nwsaom} is understood. Both commands work for a two-mode model: the simulation runs on the actors-by-events network, the table shows the number of ties (attendances) only (the columns for mutual dyads and transitive triples are missing), and \opt{generate()} saves two-mode networks (one row per actor, one column per event). A \opt{theta()} matrix of the wrong shape stops with a message that says what was expected. After a fit with \opt{groups()}, \texttt{estat simulate} is not available.

\section{Score-type tests: \texttt{estat score}}

The score-type test of Schweinberger (2012) asks, for effects that are not in the model, whether adding them would improve the fit, using only the restricted model. Fit the restricted model with the effects to be tested in the model but held fixed at the restricted value (usually 0), then call \texttt{estat score}:
\begin{lstlisting}
. nwsaom, wave1(glasgow1) wave2(glasgow2) outdegree reciprocity transtrip cycle3 theta0(-2.3 2 0 0) fix(transtrip cycle3) k0(100) k3(1000) seed(1)
. estat score

Score-type test of the restricted model (as RSiena score.Test)
Tests whether the fixed effects below differ from their fixed values.

            effect |  Fixed at     chi2  d.f. Prob>chi2 One-sided z
-------------------+-----------------------------------------------
         transtrip |    0.0000   81.920     1    0.0000       9.051
            cycle3 |    0.0000   67.902     1    0.0000       8.240
-------------------+-----------------------------------------------
        Joint test |             82.061     2    0.0000
-------------------------------------------------------------------

One-step estimates (restricted estimate plus the one-step correction)
            effect |  Estimate
-------------------+----------
         outdegree |   -2.6717
       reciprocity |    2.0383
         transtrip |    1.3303
            cycle3 |    0.2607
------------------------------
\end{lstlisting}
The joint statistic $c$ is chi-squared with as many degrees of freedom as tested effects. The one-sided variate is a normal variate with the sign changed, as in RSiena, so that it is positive when the effect is expected to be positive. The one-step estimates give the value each coefficient would approximately take if the tested effects were freed, which is a quick way to see the size of the effects before fitting the larger model. \opt{tested(\emph{effects})} restricts the test to some of the fixed effects; the others stay fixed and are treated as redundant. The test works for network effects and for the effects of a co-evolving behavior.

\paragraph{After \texttt{nwsaom multinet}.} The restricted model of a multinet fit is specified with \opt{theta0()} and \opt{fix()} of that command, with the coefficient names of the table (for example \opt{fix(net1\_x\_crprod)}); the rates of all networks and of the behavior are estimated along. The example fixes the cross-network effect \texttt{crprod} at 0 (the data are those of Section~\ref{sec:multinet}).
\begin{lstlisting}
. nwsaom multinet, net1(glasgow1 glasgow2 glasgow3) net2(advice1 advice2 advice3) eff1(outdegree reciprocity transtrip) eff2(outdegree reciprocity) cross1(crprod) cross2(from) theta0(-2.6 2.5 0.6 0 -1.9 1 0) fix(net1_x_crprod) seed(1) k3(300)
. estat score

Score-type test of the restricted model (as RSiena score.Test)
Tests whether the fixed effects below differ from their fixed values.

            effect |  Fixed at     chi2  d.f. Prob>chi2 One-sided z
-------------------+-----------------------------------------------
     net1_x_crprod |    0.0000    2.323     1    0.1275      -1.524
-------------------------------------------------------------------

One-step estimates (restricted estimate plus the one-step correction)
            effect |  Estimate
-------------------+----------
    net1_outdegree |   -2.6173
  net1_reciprocity |    2.5168
    net1_transtrip |    0.6207
     net1_x_crprod |   -0.4489
    net2_outdegree |   -1.9142
  net2_reciprocity |    1.1544
       net2_x_from |    0.2009
             rate1 |    5.5497
             rate2 |    5.0184
             rate3 |    1.8038
             rate4 |    1.7646
------------------------------
\end{lstlisting}
On the artificial advice network \texttt{crprod} has no evidence in its favour, as in the fit of Section~\ref{sec:multinet}. The rates are listed as \texttt{rate1}, \texttt{rate2} (the first network in periods 1 and 2) and \texttt{rate3}, \texttt{rate4} (the second network).

\section{Time heterogeneity: \texttt{estat timetest}}

\texttt{estat timetest} asks whether the effects of a model with at least three waves have the same size in all periods (RSiena's \texttt{sienaTimeTest}). For each effect the model is extended by deviations for the periods 2 to $P$, and the hypothesis that all deviations are 0 is tested by a score-type test, from the simulations of the fit in memory; no second fit is needed.
\begin{lstlisting}
. nwsaom, waves(glasgow1 glasgow2 glasgow3) outdegree reciprocity transtrip k0(100) k3(1000) seed(1)
. estat timetest
Test of time heterogeneity (score-type test; as RSiena sienaTimeTest)
Is the parameter of an effect the same in all periods?

            effect |      chi2  d.f.  Prob>chi2
-------------------+---------------------------
         outdegree |     0.074     1     0.7860
       reciprocity |     0.963     1     0.3264
         transtrip |     0.439     1     0.5075
-------------------+---------------------------
        Joint test |     3.567     3     0.3122
-----------------------------------------------

Period-specific deviations (one-step estimates) and separate tests
            effect |  Period   Deviation      chi2  Prob>chi2
-------------------+-----------------------------------------
         outdegree |       2     -0.3527     0.074     0.7860
       reciprocity |       2      0.6927     0.963     0.3264
         transtrip |       2      0.0779     0.439     0.5075
-------------------------------------------------------------
\end{lstlisting}
With $P$ periods there are $P-1$ deviations per effect. In this example there is no evidence that the effects change between the two periods. The tests are most informative when the deviation of an effect is the question: the deviation is reported as the one-step estimate of the difference between the period and the first period. \opt{effects(\emph{names})} restricts the test to some effects.

\paragraph{Treatment of the rates.} By default the rate parameters are treated as estimated, so that the test equals the score test of a model that adds one fixed period dummy to the fit (what RSiena's \texttt{score.Test} gives after such a fit). The option \opt{fixrates} treats the rates as known, which is what RSiena's \texttt{sienaTimeTest} does. The two agree closely except for effects whose statistic is strongly correlated with the rate, typically the outdegree effect of a network whose changes are close to random; there the default gives the larger value. In the example above they are identical to four decimals; \opt{estat timetest, fixrates} gives the same table.

\paragraph{After \texttt{nwsaom multinet}.} The test is also available for the effects of all networks and of the behavior of a multinet fit.
\begin{lstlisting}
. estat timetest
Test of time heterogeneity (score-type test; as RSiena sienaTimeTest)
Is the parameter of an effect the same in all periods?

            effect |      chi2  d.f.  Prob>chi2
-------------------+---------------------------
    net1_outdegree |     0.077     1     0.7808
  net1_reciprocity |     0.998     1     0.3177
    net1_transtrip |     0.396     1     0.5291
    net2_outdegree |     1.487     1     0.2227
  net2_reciprocity |     3.592     1     0.0580
       net2_x_from |     1.501     1     0.2205
-------------------+---------------------------
        Joint test |     7.105     6     0.3113
-----------------------------------------------
\end{lstlisting}
(Run after the multinet fit above with \texttt{crprod} fixed. The output continues with the period-specific deviations.)

\section{Relative importance of effects: \texttt{estat importance}}

\texttt{estat importance} computes the relative importance of the effects of a network model (RSiena's \texttt{sienaRI}, \texttt{interpret\_size}). At the start network of each period the contribution of every effect to every possible tie change of every actor is computed. The importance of an effect for an actor is the $L^1$ distance between the actor's choice probabilities with the fitted coefficients and with the effect's coefficient set to 0; the relative importance is its share among all effects, averaged over the actors. The command also reports $\sigma$ (the within-actor standard deviation of the effect's contributions), the semi-standardised estimate (the coefficient times $\sigma$, which compares effects on a common scale) and the certainty $R_H$ of the actors' choices.
\begin{lstlisting}
. nwsaom, waves(glasgow1 glasgow2 glasgow3) outdegree reciprocity transtrip cycle3 k0(100) k3(500) seed(1)
. estat importance
            effect | Estimate  Std. err.    Sigma  Semi-std.     RI p1     RI p2
-------------------+------------------------------------------------------------
         outdegree |  -2.7041     0.1146   0.4384    -1.1854    0.4309    0.5015
       reciprocity |   2.5066     0.2064   0.2140     0.5364    0.3763    0.3085
         transtrip |   0.6853     0.1514   0.6514     0.4464    0.1741    0.1727
            cycle3 |  -0.1348     0.2977   0.3374    -0.0455    0.0187    0.0173
-------------------+------------------------------------------------------------
     Certainty R_H |                                            0.2674    0.3548
--------------------------------------------------------------------------------
\end{lstlisting}
The command is available after fits of a directed or non-directed one-mode network or of a two-mode network, with or without one co-evolving behavior, without endowment and creation effects, \opt{structural()}, \opt{present()} or missing-data options; RSiena has the same restrictions. It needs the native plugin. All sizes agree with \texttt{sienaRI} to four decimals.

\paragraph{Non-directed and two-mode networks.} For a non-directed network the importance is that of the forcing model, as RSiena notes. For a two-mode network the actors choose among the events. With the simulated $40\times20$ data of Section~\ref{sec:twomode}:
\begin{lstlisting}
. nwsaom, waves(tm1 tm2 tm3) outdegree cycle4 nodeocov(age) nodeicov(size) seed(1) k3(300)
. estat importance
            effect | Estimate  Std. err.    Sigma  Semi-std.     RI p1     RI p2
-------------------+------------------------------------------------------------
         outdegree |  -0.8968     0.1182   0.6843    -0.6136    0.7242    0.6962
     nodeicov_size |  -0.1309     0.0748   1.2020    -0.1574    0.1678    0.1533
      nodeocov_age |   0.0139     0.0909   0.7709     0.0107    0.0136    0.0129
            cycle4 |   0.0550     0.0336   2.6018     0.1431    0.0943    0.1376
-------------------+------------------------------------------------------------
     Certainty R_H |                                            0.0771    0.0732
--------------------------------------------------------------------------------
\end{lstlisting}
\paragraph{Co-evolution.} After a co-evolution fit a second table gives the behavior effects. The options of an actor are to change the behavior by $-1$, $0$ or $+1$; the importance of a behavior effect is again the $L^1$ distance of these choice probabilities with and without the effect. The results are \texttt{r(ri\_beh)}, \texttt{r(importance\_beh)}, \texttt{r(sigma\_beh)}, \opt{r(sigma\_period\_beh)} and \texttt{r(certainty\_beh)}.
\begin{lstlisting}
. nwsaom, waves(glasgow1 glasgow2 glasgow3) outdegree reciprocity transtrip behavior(smoke1 smoke2 smoke3 : linear quadratic avalt) k0(100) k3(500) seed(1)
. estat importance
Network effects
            effect | Estimate  Std. err.    Sigma  Semi-std.     RI p1     RI p2
-------------------+------------------------------------------------------------
         outdegree |  -2.6446     0.1408   0.4384    -1.1593    0.4444    0.5151
       reciprocity |   2.4235     0.2075   0.2140     0.5186    0.3895    0.3204
         transtrip |   0.6070     0.0942   0.6514     0.3954    0.1661    0.1646
-------------------+------------------------------------------------------------
     Certainty R_H |                                            0.2604    0.3396
--------------------------------------------------------------------------------

Behavior effects (the options are changing by -1, 0, +1)
            effect | Estimate  Std. err.    Sigma  Semi-std.     RI p1     RI p2
-------------------+------------------------------------------------------------
        beh_linear |  -1.9183     0.3811   0.7517    -1.4419    0.6001    0.5875
     beh_quadratic |   2.0247     0.4437   0.6712     1.3590    0.1808    0.2095
         beh_avalt |   1.9822     1.4992   0.4392     0.8707    0.2191    0.2030
-------------------+------------------------------------------------------------
     Certainty R_H |                                            0.5814    0.5745
--------------------------------------------------------------------------------
\end{lstlisting}

\section{Several groups: \texttt{nwsaom meta}}\label{sec:meta}

This section is about the meta-analysis of separate fits. To estimate one model for all groups at once, use the option \opt{groups()} (Chapter~\ref{ch:groups}); Section~\ref{sec:groups-vs-meta} compares the two with the same data.

When the same model is fitted to several networks of the same kind, for example school classes, the group estimates can be combined, which is what RSiena's \texttt{siena08} does. \texttt{nwsaom meta} takes the stored estimates (\texttt{estimates store}) of the same model fitted to each group, with the same effects in the same order. (The name \texttt{nwsaom\_meta}, used by earlier versions, still works.) For every effect it reports the combined mean and its standard error (an iteratively reweighted least squares estimate with a random effect for the group, Snijders and Baerveldt 2003), the estimated standard deviation between groups, a heterogeneity test $Q$, the maximum likelihood estimates of the mean and the between-group standard deviation with profile-likelihood intervals, and Fisher's combination of the one-sided tests.

The following example runs with the data of the package. It makes four groups by simulating the model four times from the same start network at known coefficients (Section~\ref{sec:sim}), fits each group on its own, and combines the fits. The combined estimates should be close to the values used for the simulation ($-2.3$, 2 and 0.5).
\begin{lstlisting}
. nwwebuse glasgow, nwclear
. nwsaom simulate, wave1(glasgow1) wave2(glasgow2) outdegree reciprocity transtrip theta(-2.3 2 0.5) rates(5) nsim(4) seed(7) generate(g)
. forvalues s = 1/4 {
.     quietly nwsaom, wave1(glasgow1) wave2(g_`s'_p1) outdegree reciprocity transtrip k3(500) seed(`s')
.     estimates store grp`s'
. }
. nwsaom meta grp1 grp2 grp3 grp4
Meta-analysis of 4 fits (as RSiena siena08)
Estimates with a standard error of 10 or more are excluded;
tests of the means use a t distribution with n-1 d.f.
IWLS estimates (Snijders and Baerveldt 2003)
      effect |   n      Mean  Std. err.       t   P>|t|      SD       Q     P>Q
-------------+-----------------------------------------------------------------
   outdegree |   4   -2.2509     0.0579  -38.87   0.000  0.1158    2.59   0.460
 reciprocity |   4    1.8694     0.0737   25.37   0.000  0.1473    1.18   0.759
   transtrip |   4    0.4758     0.0114   41.81   0.000  0.0228    0.14   0.987
\end{lstlisting}
The combined means recover the coefficients, and no effect differs between the groups (\texttt{P>Q} between 0.46 and 0.99), as it should for groups simulated from the same model. In a real analysis a small \texttt{P>Q} for an effect means that the groups differ in that effect by more than their standard errors allow, and the estimated between-group standard deviation shows by how much. The full set of results is returned in \texttt{r(meta)}; the option \opt{rates} combines the rate parameters as well, and \opt{bound(\#)} excludes groups with a standard error of at least \# (default 10), as RSiena does. An effect that a group fitted with \opt{fix()} has a standard error of 0 and is left out for that group, as \texttt{siena08} leaves it out (version 1.4 and earlier counted such a group in $n$ and in the degrees of freedom of the tests of that effect, although it did not enter the weights; the estimates and the heterogeneity statistics of the other effects are unchanged). The option \opt{rates} uses the standard errors of the rates, \texttt{e(rates\_se)}, of each fit (version 1.1 took them from a row that does not exist, so that every rate was dropped without a message); the rows are named \texttt{rate\_period1}, \texttt{rate\_period2}, \ldots. The matrices \texttt{r(theta)} and \texttt{r(se)} (one row per stored estimate, one column per effect) carry the names of the stored estimates and of the effects. The results reproduce RSiena's functions to five decimals (Chapter~\ref{ch:validation}).

\section{What is not available}

After \texttt{nwsaom multinet}, \texttt{estat score} and \texttt{estat timetest} are available, but \texttt{estat gof}, \texttt{estat importance} and \texttt{estat simulate} are not. The same holds after a fit with \opt{groups()} (Chapter~\ref{ch:groups}), where \texttt{estat timetest} also reports a test per group and per period. The generalised method of moments of RSiena is not implemented (Chapter~\ref{ch:missing}). Maximum likelihood and Bayesian estimation are described in Chapter~\ref{ch:estimation}.
